\chapter{Introdução}\label{cap:introducao}

A dissipação de calor em dispositivos eletrônicos compactos motiva o
desenvolvimento de tecnologias de resfriamento eficientes. O escoamento
bifásico com mudança de fase em microcanais é uma alternativa para essa
aplicação, pois permite dissipar elevados fluxos de calor em pequenas
áreas \cite{cheng2023}. Nesse processo, o fluido refrigerante absorve calor
ao percorrer o canal e passar do estado líquido para o vapor.

Para o controle e a otimização desses sistemas, é necessário identificar o padrão de
escoamento, isto é, a configuração geométrica assumida pelas fases líquida e vapor dentro
do canal. Padrões como \textit{dispersed bubbles}, \textit{plug}, \textit{slug}, \textit{churn} e
\textit{wavy annular} apresentam comportamentos térmicos e hidrodinâmicos distintos, relevantes
para avaliar a transferência de calor e a queda de pressão \cite{cheng2023,deoliveira2021}. A identificação dos padrões por
inspeção visual de imagens de câmera de alta velocidade é subjetiva e depende
do observador \cite{dias2025knn}, o que motiva o desenvolvimento de métodos
automáticos de classificação.

No Laboratório de Estudos Térmicos e Fluidodinâmicos (LETEF) da UNISINOS,
\citeonline{deoliveira2021} caracterizam os padrões de isobutano (R600a) em microtubo
horizontal de 1\,mm pela transformada \textit{Wavelet} contínua (CWT) das séries
temporais de pressão diferencial e do sensor óptico. \citeonline{dias2025knn} e
\citeonline{dias2025cnn} apresentam classificadores baseados em imagens da seção
de visualização. Esses estudos motivam a investigação de descritores que permitem examinar quais
características das imagens se repetem entre ensaios de um mesmo padrão.

Neste trabalho, compararam-se as escalas das variações de intensidade entre
ensaios a partir de panoramas. Cada panorama foi construído pela junção de
partes de quadros consecutivos da câmera, formando uma imagem alongada.
Dessa imagem, extraiu-se um sinal unidimensional de intensidade em função da
posição em pixels. Como reuniu-se registros de instantes distintos, o panorama
não correspondeu a uma fotografia instantânea do escoamento. Sua construção
e a extração do sinal foram detalhadas na \autoref{sec:met-extracao}.

\citeonline{schuck2019} apresentam uma formulação alternativa para o cálculo da
coerência \textit{Wavelet} (WCO) em
que a transformada é aplicada às funções de autocorrelação e correlação cruzada dos sinais
em vez de ser aplicada diretamente a eles. Essa abordagem elimina a necessidade de
suavização espectral exigida pela definição clássica \cite{torrence1998,grinsted2004}. A questão
investigada é se a forma dos perfis de potência dos sinais reconstruídos apresenta
maior semelhança entre ensaios do mesmo padrão do que entre padrões distintos.
Os mapas de coerência complementam essa comparação pela leitura das associações
entre sinais em função da escala e do atraso e por resumos numéricos de cada par.
A avaliação é descritiva e não pressupõe que
correlação elevada demonstre identidade do padrão.

Como complemento, examinaram-se a fração de vazio estimada a partir de máscaras de
gás e a densidade de transições dessas máscaras ao longo dos panoramas. A primeira
indicou a ocupação estimada por gás, enquanto a segunda registrou a alternância espacial
entre regiões segmentadas. Essas informações auxiliaram a interpretação dos perfis
\textit{Wavelet}, que descrevem as escalas das variações de intensidade. A combinação de
informações espectrais e de fração de vazio encontra um antecedente em
\citeonline{nguyen2010}. Neste trabalho, estimou-se a ocupação por gás a partir
das regiões identificadas como gás nas imagens, supondo seções gasosas circulares.
A estimativa dependeu, portanto, da identificação dessas regiões e da forma
geométrica assumida. Não se realizou uma calibração por comparação com uma
medição independente do volume real de gás.

% =========================================================
\section{Objetivo Geral}
% =========================================================

O objetivo geral desse trabalho é analisar as semelhanças e diferenças entre
ensaios de um mesmo padrão de escoamento bifásico e de padrões distintos no
arranjo experimental do LETEF, a partir da representação \textit{Wavelet} dos
sinais espaciais extraídos dos panoramas. A análise se concentra na comparação
dos perfis de potência pelo coeficiente de Pearson $r$, na interpretação dos
mapas e nos resumos de coerência, com apoio das medidas ópticas de fração de vazio
estimada e alternância espacial das fases.

% =========================================================
\section{Objetivos Específicos}
% =========================================================

\begin{itemize}
  \item Interpretar os mapas de potência, espectro cruzado e coerência
  \textit{Wavelet} em exemplos de pares do mesmo padrão e de padrões distintos,
  considerando o comprimento de onda e o atraso em pixels.
  \item Resumir os valores de coerência \textit{Wavelet} nos mapas de cada par de
  ensaios, considerando todas as escalas e atrasos analisados, utilizando o
  mínimo e o máximo para identificar os extremos e a média e a mediana para
  descrever os valores centrais.
  \item Comparar a semelhança da forma dos perfis espectrais dentro de cada padrão
  e entre padrões distintos pelo coeficiente de Pearson $r$, considerando
  mínimos, médias, medianas, máximos e a sobreposição das distribuições.
  \item Examinar a complementaridade entre os resultados \textit{Wavelet}, a estimativa
  óptica de fração de vazio e a densidade espacial de transições das máscaras
  de gás, comparando as medidas dos mesmos ensaios.
  \item Delimitar a interpretação dos resultados a partir da qualidade dos
  panoramas, das condições de aquisição e das hipóteses adotadas no
  pré-processamento e nas análises.
\end{itemize}
